% H9-T18: positive-characteristic threefold-cover gate.
% Run JAC1964-20261121; /root/h9_height; GPT-6 Astra Pro.

\subsection*{H9-T18. The cover gate reduces to the variable middle form of one incidence fourfold}

Let $A/k$ be an abelian fourfold over an algebraically closed field,
with ample class $h$, and let a genuine Lefschetz pencil have blowup
$\widetilde A\to\mathbf P^1$. Let $\pi:C\to\mathbf P^1$ be a finite
separable cover unramified above the critical values. Put
\[
 M=A\times C,\qquad
 \mathcal X=\widetilde A\times_{\mathbf P^1}C,
 \qquad k:\mathcal X\hookrightarrow M.
\]
Then $\mathcal X$ is a smooth projective ample divisor of class
\[
 \lambda=h+b e,\qquad e=[\text{point of }C],\qquad b=\deg\pi>0.
\]
Let $Y/k(C)$ be its generic fibre and let
$\mathcal L=\lambda|_{\mathcal X}$, whose generic restriction is
$h_Y$. Every homologically trivial cycle on $Y$ is admissible by
H9-T12. Throughout, ``algebraic classes'' means images of rational
Chow groups in one fixed invertible-prime cohomology.

The products of a curve and an abelian variety have algebraic
Kunneth and Lefschetz operators. This follows either from the
classical product stability of $B$, or from JAC-P16, since the
curve motive is a summand of its Jacobian motive together with
Tate motives. Thus all inverse operators used below are algebraic.

\paragraph{The fixed primitive form is positive by032, with all factors explicit.}
An algebraic class $a\in H^4(M)$ has a Kunneth decomposition
\[
                 a=\beta+\eta+\gamma e,
 \quad\eta\in H^3(A)\otimes H^1(C),\quad\gamma\in H^2(A).
\]
The curve Kunneth projectors and the abelian Lefschetz projectors
are algebraic, so each displayed part is algebraic. Suppose
$\lambda^2a=0$. The component equations are
\[
 h^2\beta=0,\qquad h^2\eta=0,
       \qquad h^2\gamma+2b h\beta=0.
\]
The Lefschetz decomposition on $A$ and the first equation give
$\beta=\beta_0+h\alpha$, where $h\beta_0=0$ and $h^3\alpha=0$.
There is no scalar $h^2$ term. Since
$h^2:H^2(A)\to H^6(A)$ is an isomorphism, the last equation gives
$\gamma=-2b\alpha$. Consequently
\[
 a=\beta_0+(h\alpha-2b\alpha e)+\eta,
\]
and direct multiplication gives the orthogonal norm identity
\[
 \int_M a^2\lambda
   =b\int_A\beta_0^2
       -3b\int_Ah^2\alpha^2
       +\int_{C\times A}\eta^2h.                          \tag{H9.86}
\]
The mixed terms vanish by degree or by $h\beta_0=0$.

The first two terms are positive for nonzero components by032,
Corollary 8.4, in codimensions two and one respectively. For the
last term, write $g$ for the genus of $C$, distinct from the
positive degree $b$. If $g=0$, $\eta=0$. If $g>0$, the algebraic
inverse to Abel pullback lifts $\eta$ to an algebraic class
\[
 w\in H^1(J(C))\otimes H^3(A),\qquad h^2w=0.
\]
For the principal Jacobian polarization $\theta$, put $H=\theta+h$.
Then $H^{g+1}w=0$: the terms containing $h^2$ vanish, and every
remaining term contains $\theta^gH^1(J(C))=0$. This is the
primitive exponent for a codimension-two class on the
$(g+4)$-dimensional abelian variety $J(C)\times A$. Degree and
the Abel curve-class formula give
\[
 \int_{J(C)\times A}w^2H^g
       =g!\int_{C\times A}\eta^2h.                        \tag{H9.87}
\]
Thus the third term also has the required strict positive sign by
032. Equation (H9.86) proves positive definiteness on the
primitive algebraic $H^4(M)$, in every characteristic. It does
not assume that arbitrary motivic summands inherit their own
geometric Hodge-standard forms.

\paragraph{The remaining variable subspace is explicit.}
Hard Lefschetz on $M$ makes
$L_\lambda:H^4(M)\to H^6(M)$ an isomorphism. Define
\[
 r=L_\lambda^{-1}k_*,\qquad
 V_{\mathcal X}=\ker(k_*:H^4(\mathcal X)\to H^6(M)(1)).
\]
These operators give an algebraic orthogonal decomposition
\[
 H^4(\mathcal X)=k^*H^4(M)\ \mathbin\perp\ V_{\mathcal X}.
                                                               \tag{H9.88}
\]
Indeed $rk^*=1$, and projection formula proves orthogonality.
The Gysin map $k_*:H^6(\mathcal X)\to H^8(M)(1)$ is injective,
by dual weak Lefschetz in degree two. Hence every variable class
is primitive: $k_*(\mathcal L v)=\lambda k_*v=0$ forces
$\mathcal L v=0$. An ambient class $k^*a$ is primitive exactly
when $\lambda^2a=0$. Its norm is (H9.86). In particular the only
uncontrolled primitive algebraic form in (H9.88) is
\[
        v\longmapsto\int_{\mathcal X}v^2,
            \qquad v\in\operatorname{Alg}^2(\mathcal X)
                                  \cap V_{\mathcal X}.      \tag{H9.89}
\]
The fact that $M$ is of abelian type gives no Chow-motive splitting
of this new variable part into abelian motives.

\paragraph{Zhang's primitive correction can be made algebraic in this setting.}
Let $z$ be an admissible extension on $\mathcal X$ of a
homologically trivial generic cycle with
$\operatorname{alb}(h_Yz_Y)=0$. Consider
\[
         \ell(D)=\int_{\mathcal X}z\mathcal L D,
                                      \qquad D\in\operatorname{Pic}(\mathcal X).
\]
Vertical orthogonality makes this depend only on $D_Y$.
Kahn Theorem 6.2 makes it zero on $\operatorname{Pic}^0(Y)$,
by the Albanese hypothesis. Thus it is a functional on
$\operatorname{NS}(Y)_{\mathbf Q}$. The nondegenerate numerical
pairing between divisors and one-cycles supplies a rational
one-cycle $W_Y$ such that
$\ell(D)=\deg_Y(D_YW_Y)$.

Here $Y$ is a smooth hyperplane threefold in $A_{k(C)}$.
The inverse to $h_Y:H^2(Y)\to H^4(Y)$ is algebraic: if
$j:Y\hookrightarrow A$, it is
\[
                     j^*(L_h^2)^{-1}j_*.
\]
Indeed $j_*:H^4(Y)\to H^6(A)$ is an isomorphism by dual weak
Lefschetz, and $j_*h_Yj^*a=h^2a$. Apply this inverse to $W_Y$
to obtain a rational divisor $U_Y$ with
$h_YU_Y=W_Y$ in cohomology, hence numerically. Close $U_Y$ to
a horizontal divisor $U$ on $\mathcal X$, and let $u$ be its
fibre over a point of $C$. Then
\[
        \delta=z-u,\qquad
        \int_{\mathcal X}\delta\mathcal L D=0
                         \quad\text{for every divisor }D.  \tag{H9.90}
\]
Moreover $u^2=0$ and $z\cdot u=0$, so $\delta^2=z^2$ in degree.

We can strengthen (H9.90) to actual cohomological primitivity
without using a characteristic-zero Hodge decomposition. The
class
\[
       c=(L_\lambda^3)^{-1}\bigl(\lambda k_*\delta\bigr)
                                        \in H^2(M)
\]
is algebraic, because the inverse on $M$ is algebraic. For every
algebraic divisor class $d$ on $M$, (H9.90) gives
\[
                  \int_M c\,d\,\lambda^3=0.
\]
The ordinary divisor Hodge index theorem, together with numerical
versus homological equivalence for divisors, makes this form
nondegenerate on divisor classes. Hence $c=0$. The injectivity of
$k_*:H^6(\mathcal X)\to H^8(M)$ then gives
$\mathcal L\delta=0$.
This proves, in this particular incidence model and in every
characteristic, the primitive correction needed by Zhang's
argument. The final Hodge-index sign on (H9.89) is the remaining
input, not an additional hidden algebraicity hypothesis in the
correction step.

Splitting $[\delta]=k^*a+v$ by (H9.88) now gives
\[
 H_Y(z_Y,z_Y)
      =\int_M a^2\lambda+\int_{\mathcal X}v^2,
       \qquad \lambda^2a=0.                               \tag{H9.91}
\]
The first term is positive definite by (H9.86)--(H9.87).
Therefore semipositivity of (H9.89) suffices for semipositivity
on the full admissible Albanese-primitive generic group.

\paragraph{The reduction is an equivalence on this unramified-cover model.}
Conversely let an algebraic variable class $v$ be represented by
$\gamma\in\operatorname{CH}^2(\mathcal X)_{\mathbf Q}$.
Its restriction to the generic fibre is homologically trivial.
To check this, pair it with $H^2(Y)=j^*H^2(A)$: every such
pairing equals
$\int_Mk_*v\,\operatorname{pr}_C^*[\mathrm{point}]
                          \operatorname{pr}_A^*a=0$.
Poincare duality gives the assertion. H9-T12 then makes $\gamma$
admissible on this model.

For every $D\in\operatorname{Pic}^0(Y)_{\mathbf Q}$, choose a
divisor closure on $\mathcal X$. Weak Lefschetz identifies $H^2(\mathcal X)$ with
$k^*H^2(M)$, so $k_*v=0$ gives
\[
                         H_Y(\gamma_Y,h_YD)=0.              \tag{H9.92}
\]
The corresponding global intersection vanishes for arbitrary
divisor closures, but the height identity is used only for these
algebraically trivial $D$, for which $h_YD$ lies in its domain.
In the rational decomposition of H9-T17, write
$\gamma_Y=z_0+h_YD_\gamma$ with
$\operatorname{alb}(h_Yz_0)=0$. Orthogonality and (H9.92) make
the negative Picard height of $D_\gamma$ zero; Kahn Theorem 6.6
therefore puts $D_\gamma$ in the trace subgroup.

For these hyperplane fibres, Picard restriction from the constant
$A$ is an isogeny, and the isogeny
$D\mapsto\operatorname{alb}(h_Y^2D)$ is constant: after the
Albanese identification its formula is intersection with $h^3$
on $A$. Thus, after rational normalization, $D_\gamma$ is the
restriction of a constant algebraically trivial divisor on $A$.
Its global product with $\mathcal L$ has zero cohomology class.
Subtracting that product from $\gamma$ leaves its variable class
$v$ unchanged and gives the Albanese-primitive generic cycle
$z_0$, with
\[
                          H_Y(z_0,z_0)=\int_{\mathcal X}v^2.
\]
Consequently positivity on the full Albanese-primitive generic
group implies semipositivity of (H9.89). This proves the claimed
equivalence. If (H9.89) is positive definite on algebraic
cohomology, (H9.91) also identifies the primitive height radical
with the restrictions of globally homologically trivial cycles,
intersected with the Albanese-primitive domain.

\paragraph{What remains to be proved.}
Accepted032 settles all three fixed terms in (H9.86). It does
not settle (H9.89) merely because the fourfold is an ample divisor
in a product of a curve and an abelian variety. For the original
rational pencil base, the blowup decomposition replaces this
variable form by the surface-axis form, which is why H9-T8 closes
that case. After a separable cover the new variable fourfold
classes are exactly the additional issue identified here.

A possible next approach is monodromy while moving the branch
parameters of the cover. That is a family of incidence divisors
formed from two pencil sections. It is generally a rank-two
subsystem of the full hypersurface linear system on $A\times C$.
A generic-hypersurface full-orthogonal-monodromy theorem cannot
be applied to that subsystem without a new proof. The present
result is an exact reduction with a proved positive fixed part,
not a claim that the new variable form vanishes or is positive.

The sources are the previously source-checked Kahn Theorems 4.14,
5.6, 6.2 and 6.6, the ordinary all-characteristic divisor index
and divisor-equivalence theorems, and accepted032 Corollary8.4.
The new calculations and the two directions of the reduction are
proved above. No general Hodge-standard claim for arbitrary
abelian-type varieties is inferred.
